% !TeX root = ../main.tex


\newlecture[Independence, Random Variables and Distributions][Sep 15, 2026]

\deflabel{2.7}{Independence}
\begin{quote}
    Events $E_1, \dots, E_n \subseteq \Omega$ are independent with respect to a probability measure $P$, if for each $\mathcal{I} \subseteq \underset{\text{any subset of} \left\{1, \dots, n\right\}}{[n]}$
    \[P\left(\bigcap_{i\in\mathcal{I}}E_i\right) = \prod_{i\in\mathcal{I}}P(E_i)\]
\end{quote}
e.g. if $\mathcal{I} = \left\{1,2,3\right\}$ a subset of $\left\{1,\dots,n\right\}$
\[P(E_1 \bigcap E_2 \bigcap E_3) = P(E_1)P(E_2)P(E_3)\]

\begin{proposition}[][2.7.1]
    If $E_1, \dots E_n$ are independent, then $E_1^c, \dots, E_n^c$ are also independent. Proof as an exercise.
\end{proposition}

\begin{proposition}[][2.7.2]
    An infinite collection of events $\left\{E_\alpha :\alpha \in \mathcal{I}\right\}$ are independent. Then for a finite subset $\mathcal{J} \subseteq \mathcal{I}$, the events $\left\{E_\beta : \beta \in \mathcal{J}\right\}$ are independent.
\end{proposition}

\subsection{Random Variables}
\begin{definition}[Random Variables][3.1]
    Given a sample space $\Omega$, a function $X: \Omega \rightarrow \mathbb{R}$ is a random variable.
\end{definition}

\begin{example}[Coin Flipping][3.1.1]
    \[\Omega = \left\{H, T\right\}\]
    \[\downarrow \text{R.V. } X: \quad X(H)\rightarrow 1, \quad X(T)\rightarrow 0 \]
    \[\mathbb{R}\]
\end{example}

\begin{theorem}[Distribution][3.2]
    A random variable $X$ and probability measure $P$ induce a probability measure $\mu$ on $R$ defined by
    \[\mu(A) = p(X \in A) = p(\left\{\omega\in \Omega: X(\omega)\in A\right\})\]
    for any $A \subseteq R$
    \aside{$\mu$ is a measure on the real number set: $\mu(\{0,1\})$ and $\mu([0,1))$ each return a certain value.}
    \\ $\mu(A)$ takes a real number set $A$ as input, and equals $P$ of the \aside{pre-image of $A$ w.r.t.\ r.v.\ $X$}, i.e.\ $P(\Omega_A) = \mu(A)$ where $\Omega_A = \left\{\omega\in\Omega: X(\omega)\in A\right\}$.
\end{theorem}

\begin{example}[][3.2.1]
    Define the following
    \[\Omega = \left\{ \text{\textcolor{blue}{blue}, \textcolor{red}{red}, \textcolor{green}{green}, \textcolor{yellow}{yellow}} \right\} \]
    \[X: \begin{cases}
        \textcolor{blue}{\text{blue} \rightarrow 1} \\
        \textcolor{red}{\text{red} \rightarrow 3} \\
        \textcolor{green}{\text{green} \rightarrow 4} \\
        \textcolor{yellow}{\text{yellow} \rightarrow 6}
    \end{cases}\]

    \[P: \{\text{blue}\}\mapsto\tfrac{1}{5},\quad \{\text{red}\}\mapsto\tfrac{1}{5},\quad \{\text{green}\}\mapsto\tfrac{1}{4},\quad \{\text{yellow}\}\mapsto 1-\tfrac{1}{5}-\tfrac{1}{5}-\tfrac{1}{4}=\tfrac{7}{20}\]
    \[P(\{\text{blue, red}\}) = \tfrac{1}{5}+\tfrac{1}{5} = \tfrac{2}{5}\]
    \[A = [1,3] \]
    \begin{enumerate}
        \item $\left\{w: w \in \Omega, X(w) \in A \right\} = \left\{\text{\textcolor{blue}{blue} ,\textcolor{red}{red}}\right\}$
        \item $P(\left\{\text{\textcolor{blue}{blue} ,\textcolor{red}{red}}\right\}) = \frac{2}{5}$, so $\mu(A) = \frac{2}{5}$
    \end{enumerate}
    in summary
    \[\mu([1,3]) = \frac{2}{5}\]
\end{example}

\textbf{Claim:} $\mu$ is a probability measure on $\mathbb{R}$ (for any real number set $B$).
\\proof:
\begin{enumerate}
    \item $\mu(B) \in [0,1]$ is trivial, since $P$ is a probability measure.
    \item Since $X(\omega) \in R $ for all $\omega \in \Omega$.
    \[\mu(R) = p(\underset{\Omega}{ \left\{\omega : \omega \in \Omega, X(\omega)\in R\right\} }) = p(\Omega) = 1\]
    \item Consider disjoint sets $A_1, A_2, \dots \subseteq R$ and define $E_i = \left\{\omega: X(\omega) \in A_i\right\}$
    \\(e.g. $E_1 = \left\{\omega: X(\omega)\in A_1\right\}$, $E_2 = \left\{\omega: X(\omega)\in A_2\right\}$)
    \\Since $X$ is a function, it is impossible for $X(\omega) = a$ and $X(\omega) = b$ when $a\neq b$, so the $E_i$'s are also disjoint. Then
    \begin{align*}
        \mu\left(\bigcup_{i=1}^{\infty}A_i\right) &= p\left(X(\omega) \in \bigcup_{i=1}^{\infty}A_i\right) \\
        &= p\left(\bigcup_{i=1}^{\infty}E_i\right) \quad \text{we know } E_i\text{'s are disjoint} \\
        &= \sum_{i=1}^{\infty}p(E_i) \\
        &= \sum_{i=1}^{\infty}\mu(A_i)
    \end{align*}
\end{enumerate}
\hfill $\square$

\begin{definition}[The Cumulative Distribution Function][3.3]
    A random variable $X$ and measure $p$ generate the cumulative distribution function (CDF) defined by
    \[F(x) = \mu((-\infty, x]) = p(X \leq x)\]
    where $p(X\leq x)$ is a short-hand notation for
    \[p(\left\{\omega: \omega \in \Omega, X(\omega) \in (-\infty, x]\right\})\]
\end{definition}

\begin{theorem}[][3.3.1]
    A distribution function $F$ for a random variable $X$ uniquely defines the measure $\mu$.
    \\Proof is outside the scope of this class.
\end{theorem}

\begin{example}[Uniform Distribution][3.3.2]
    $P$ is uniform on $[0,1]$ and $X(\omega)=\omega$
    \\For $x\in[0,1]$
    \begin{align*}
        F(x) &= \mu((-\infty, x]) & \leftarrow& \text{by definition of CDF} \\
        &= p(X \leq x) & \leftarrow& \text{by Thm 3.2 (definition of distribution)} \\
        &= p\left(\left\{\omega :X(\omega) \in (-\infty, x]\right\}\right) \\
        &= p((-\infty, x]) & \leftarrow& \text{by }X(\omega) = \omega \\
        &= p([0,x]) & \leftarrow& \text{because $P$ is a uniform/Lebesgue measure} \\
        &= x
    \end{align*}
\end{example}

\begin{theorem}[Properties of Distribution Function][3.4]
    A distribution function satisfies the following properties:
    \begin{enumerate}[label = \roman*)]
        \item For all $x \leq y \in R, F(x) \leq F(y)$
        \item $\lim_{x \to -\infty} F(x) = 0 $ and $\lim_{x \to +\infty} F(x) = 1 $
        \item $F$ is right-continuous
        proof:
        \begin{enumerate}[label = \roman*)]
            \item $x\leq y \in R$
            \[F(x) = \mu((-\infty, x]), F(y) = \mu((-\infty, y])\]
            Since $x\leq y$, \quad $(-\infty, x] \subseteq (-\infty, y]$
            \[\mu((-\infty, x])\leq \mu((-\infty, y])\]
            "bigger set, bigger probability"
            \[(A \leq B, \quad p(A) \leq p(B) \quad / \quad \mu(A)\leq \mu(B))\]
            \item Define $A_n = \left\{X\leq -n\right\}, \quad n\geq 0$
            \\Then $A_1 \supseteq A_2 \supseteq A_3 \supseteq \dots $
            \begin{align*}
                \bigcap_{n=1}^{\infty} A_n &= \phi \\
                \lim_{n \to \infty} F(-n) &= \lim_{n \to \infty} p(A_n) \\
                &= p\left(\bigcap_{n=1}^{\infty}A_n\right) \\
                &= p(\phi) = 0
            \end{align*}
            Define $B_n = \left\{X\leq n\right\}$,
            \\then $B_1 \subseteq B_2 \subseteq \dots$
            \[\bigcup_{n=1}^{\infty} B_n = \Omega\]
            \begin{align*}
                \lim_{n \to \infty} F(n) &= \lim_{n \to \infty} p(B_n) \\
                &= p\left(\bigcup_{n=1}^{\infty}B_n\right) \\
                &= p(\Omega) = 1
            \end{align*}
            \item Fix $x\in R$ and let $h_n \downarrow 0 \quad (h_n \geq 0 \quad h_n \rightarrow 0, h_1 \geq h_2 \geq \dots)$
            \\Define $C_n = \left\{X\leq x+h_n\right\}$
            \\then $C_1 \supseteq C_2 \supseteq \dots$
            \[\lim_{n \to \infty} C_n = \bigcap_{n=1}^{\infty} C_n = \left\{X\leq x\right\} \]
            Therefore
            \[\lim_{n \to \infty} F(x+h_n) = \lim_{n \to \infty} p(C_n) = p(\bigcap_{n=1}^{\infty} C_n) = p(X \leq x) = F(x) \]
            Hence
            \[\lim_{h \downarrow 0}F(x+h) = F(x)  \]
            So $F$ is right-continuous.
        \end{enumerate}
    \end{enumerate}
\end{theorem}

\begin{definition}[The Inverse CDF][3.5]
    The inverse CDF of a r.v. $X$ is defined by
    \[F^{-1}(y) = \sup\left\{x:F(x) < y\right\}\]
\end{definition}

\begin{theorem}[][3.5.1]
    If $F$ satisfies properties i) to iii) of Thm 3.4, then such an $F$ is the CDF of some r.v. 
\end{theorem}

\begin{theorem}[][3.6]
    If $F$ is a CDF, it can only have countably many discontinuities.
\end{theorem}
